\documentclass[10pt,a4paper]{amsart}
\usepackage[margin=19mm]{geometry}
\usepackage{amsmath,amssymb,mathtools}
\usepackage[scr=rsfs,cal=euler]{mathalfa}
\usepackage{xfrac}
\usepackage[unicode,hidelinks,bookmarks=false]{hyperref}
\newcommand{\ie}{i.e.\ }
\newcommand{\cf}{cf.\ }
\newcommand{\resp}{respectively\ }
\newcommand{\Subsection}{\subsection}
\providecommand{\nobreakdash}{\nobreak\hbox{-}}
\setlength{\emergencystretch}{2em}
\multlinegap0pt
\usepackage{bm}
\usepackage{bbm}
\usepackage{dsfont}
\usepackage{xspace}
\usepackage{cancel}
\usepackage{mathtools}
\usepackage{amsthm}
\makeatletter
\@ifundefined{thm}{\newtheorem{thm}{Theorem}}{}
\@ifundefined{lemma}{\newtheorem{lemma}[thm]{Lemma}}{}
\makeatother
\newcommand{\R}{{\mathbb R}}		%Real numbers
\newcommand{\adj}{{precisely }} % \newcommand{\pre}{precisely} will not work
\newcommand{\cn}{\colon}                %ColoN
\newcommand{\la}{\langle}	   	%Left Angle
\newcommand{\ra}{\rangle}          	%Right Angle
\title[Continuous nowhere differentiable multivariate functions]{Continuous nowhere differentiable multivariate functions}
\author{Maria Girardi, Ralph Howard}
\date{Source: arXiv:2510.13061v1 (CC BY 4.0)}
\begin{document}
\maketitle

\noindent\emph{Source and licence.} Continuous nowhere differentiable multivariate functions. Version arXiv:2510.13061v1 (CC BY 4.0), distributed under the Creative Commons Attribution 4.0 International licence, as recorded in the arXiv licence metadata for this exact version and shown on the abstract page https://arxiv.org/abs/2510.13061v1. Full rights notice: licenses/arxiv-2510.13061-CC-BY-4.0.txt.\par\smallskip

\noindent\textit{Editorial scope.} Counted unit: the Lemma whose source label is \texttt{lem:not\_dense} in the article (source0.tex lines 706--709, source label \texttt{lem:not\_dense}), delivered together with the article's own complete original proof of it (source0.tex lines 712--789, one \texttt{proof} environment of 78 lines). No mathematics is written, repaired or re-derived: statement and proof are the article's own text line for line. Required context retained uncounted: thm:main\_example (lines 537--559). Every definition, display and label of those lines is retained unchanged. Omitted from this package and not redistributed: 1--536, 560--705, 710--711, 790--1232. Nothing inside a retained excerpt is omitted or altered: every retained body is delivered exactly as the article's lines are written, so no omission receipt of any kind accompanies this package. Citations: the retained excerpts carry no citation command at all, so no bibliography entry of the article is transcribed and no reference is redistributed. Reference adaptation: none was needed and none was made; every reference inside the delivered unit resolves inside it, so no unresolved reference or citation is produced. Printed slips kept literal: no printed word, symbol, hypothesis or number of the counted statement or of its proof was altered. Layout and class: the article's own class is \texttt{amsart} with packages including amsmath, amssymb, mathtools, graphicx, overpic, enumerate, verbatim, wrapfig, caption, hyperref, xcolor; the wrapper is the lane's pinned \texttt{amsart} editorial wrapper at 10pt on A4, which restates the theorem-like environments the retained text needs and adds the article's own macro definitions that the retained text needs (\texttt{\textbackslash R}, \texttt{\textbackslash adj}, \texttt{\textbackslash cn}, \texttt{\textbackslash la}, \texttt{\textbackslash ra}), with extra wrapper packages (bm, bbm, dsfont, xspace, cancel, mathtools). The delivered numbering is the unit's own. Assets: no retained span contains a figure environment, a graphics-inclusion command, a PDF-page inclusion, a table or a TikZ picture; no image file is copied into this package, the package declares no asset dependency and no image file is redistributed with it. Licence: version arXiv:2510.13061v1 of this article is distributed under the Creative Commons Attribution 4.0 International licence, as recorded in the arXiv licence metadata for this exact version and shown on the abstract page https://arxiv.org/abs/2510.13061v1. A full rights notice is delivered at licenses/arxiv-2510.13061-CC-BY-4.0.txt. Characters: every retained excerpt is pure ASCII, so the delivered engine, XeLaTeX with its default Latin Modern text font, reports no missing character.
\begin{thm}\label{thm:main_example}
  Let $\gamma\in(0,1]$ and 
  $\alpha_1, \alpha_2 ,\ldots, \alpha_d$ be distinct 
  numbers in the interval $(\tfrac{1}{1+\gamma},1)$.
  For each
  $j \in \{1,2, \ldots, d\}$, let the function 
  $f_j\cn\R\to\R$ be  \adj $C^{0,\alpha_j}$
  at each point in $\R$. 
Define  $f\cn \R^d \to \R$ by 
$$
f(x_1, x_2 ,\ldots,x_d) = \sum_{j=1}^d f_j(x_j). 
$$
Then $f$ is continuous and, 
for any $C^{1,\gamma}$ test curve $c: [a,b]\to\R^d$,
the composite $f\circ c$ is nowhere differentiable on $(a,b)$.


Moreover,
for any $C^{1,\gamma}$ test curve $c\cn [a,b]\to \R^d$
and any $s_0\in (a,b)$, there
is an $\alpha \in \{ \alpha_1, \alpha_2 ,\ldots, \alpha_d\}$ such
that $f\circ c $ is \adj $C^{0,\alpha}$ at $s_0$. 
\end{thm}

\begin{lemma}\label{lem:not_dense}
  The set $\mathcal F_n^\gamma (U)$ is a closed nowhere dense
  subset of $C(\overline U)$.
\end{lemma}

\begin{proof}
Let $\gamma \in (0,1]$. 
Let $\la f_\ell\ra_{\ell =1}^\infty$ be a sequence in $\mathcal F_n^\gamma (U)$
with  $f_\ell \to f$ uniformly  on~$\overline U$
for some   $f\in C(\overline U)$.
To show $\mathcal F_n^\gamma (U)$ is closed it is enough to show
$f\in\mathcal F_n^\gamma (U)$. 
For each $\ell$ there is
 $c_\ell \in \mathcal C_n^\gamma (U)$
with 
$$
| f_\ell(c_\ell(s)) - f_\ell(c_\ell(0))|\le n |s|
$$
for all $s\in [-1/n,1/n]$.
Since  $\mathcal C_n^\gamma (U)$ is compact, 
passing to a subsequence, we  assume there
is a $c\in \mathcal C_n^\gamma (U)$  such that
$c_\ell \to c$ in the $C^1$ topology.
Then
$$
\lim_{\ell \to \infty} f_\ell(c_\ell(s)) = f(c(s))
$$
for each  $s\in [-1/n,1/n]$ since
$f$ is continuous, $c_\ell(s)\to c(s)$, and 
\begin{align*}
  |f_\ell(c_\ell(s))- f(c(s))|&\le |f_\ell(c_\ell(s)) - f(c_\ell(s))| 
        + |f(c_\ell(s)) - f(c(s))|\\
    &\le \| f_\ell - f\|_{L^\infty} + |f(c_\ell(s)) - f(c(s))|\\
    & \stackrel{\ell\to \infty}{\longrightarrow} 0.
\end{align*}
Therefore
$$
|f(c(s)) - f(c(0))| = \lim_{\ell\to \infty} |f_\ell(c_\ell(s))-f_\ell(c_\ell(0))|
\le n |s|
$$
for all $s\in [-1/n,1/n]$.  
Thus $f\in \mathcal F_n^\gamma (U)$.
So  $\mathcal F_n^\gamma (U)$ is closed.


It remains to show $\mathcal F_n^\gamma (U)$ is nowhere dense. 
Toward a contradiction assume this is not the case.
Then the closed set $\mathcal F_n^\gamma (U)$ has nonempty interior. 
Thus $\mathcal F_n^\gamma (U)$ contains 
a nonempty  open subset $\mathcal V$ of $C(\overline U)$.
As  the  smooth functions are dense in $C(\overline U)$
there is a $f_0\in \mathcal V$ that is~$C^\infty$.
Let $f$ be the function constructed in  Theorem~\ref{thm:main_example}.


For a small enough positive $\delta$
the sum $f_\delta:=f_0+\delta f$ 
is in the open set~$\mathcal V$ and
so $f_\delta \in  \mathcal F_n^\gamma (U)$.
%  
By the definition of $ \mathcal F_n^\gamma (U)$
there is~$c\in \mathcal C_n^\gamma  (U)$ so that
the Lipschitz inequality
\begin{equation*}
    |f_\delta(c(s)) - f_\delta(c(0))| \le n |s|
\end{equation*}
holds for each  $s\in[-1/n,1/n]$.
%
However, by Theorem~\ref{thm:main_example}, $f\circ c$ is \adj~$C^{0,\alpha}$
at $0$ for some $\alpha\in (0,1)$
since the curve $c$ is $C^{1,\gamma}$.
As~$f_0$ is smooth, $f_0\circ c$ is~$C^1$.
%
Thus the sum $f_\delta\circ c = f_0\circ c +  \delta ( f\circ c) $
is \adj~$C^{0,\alpha}$ at~$0$ where $\alpha<1$.
Hence 
\begin{equation*} 
   \limsup_{s\to 0}  \frac{|f_\delta(c(s)) -f_\delta(c((0))|}{|s-0|} = \infty, 
  \end{equation*} 
contradicting  the previous 
inequality.
So  $\mathcal F_n^\gamma (U)$ is nowhere dense.
\end{proof}  

\end{document}
